\section{Architekturentscheidungen (ADRs)}

Architecture Decision Records (ADRs)\footnote{%
  \textit{Architecture Decision Record (ADR)}: Ein ADR ist ein kurzes, strukturiertes Dokument,
  das eine einzelne, architekturrelevante Entscheidung samt Kontext, betrachteten Alternativen
  und Begründung festhält. Das Format geht auf Michael Nygard zurück
  (\url{https://cognitect.com/blog/2011/11/15/documenting-architecture-decisions}).}
dokumentieren wichtige Architekturentscheidungen mit Kontext, Alternativen und Konsequenzen.
Jeder ADR ist unveränderlich – er wird nicht überschrieben, sondern durch einen neuen ADR
abgelöst oder ergänzt. Die nachfolgenden ADRs decken alle in diesem Dokument referenzierten
Entscheidungen vollständig ab.

% ── ADR-001 ──────────────────────────────────────────────────
\subsection{ADR-001: Microservices entlang der Bounded Contexts}

\begin{tabularx}{\textwidth}{L{3.5cm} X}
  \rowcolor{tableheader}
  \color{white}\textbf{Feld} & \color{white}\textbf{Inhalt} \\
  Status & Akzeptiert -- 2026-01-20 \\
  \rowcolor{tablegrey}
  Qualitätsziele & Maintainability (primär), Reliability-Availability, Performance Efficiency \\
  Entscheidungstreiber & Unabhängige Deploybarkeit der Teams; unterschiedliche Skalierungsbedarfe je BC; langfristige Wartbarkeit durch klare fachliche Grenzen (DDD). \\
  \rowcolor{tablegrey}
  Kontext & 5~klar abgrenzbare fachliche Domänen mit unterschiedlichen Änderungsraten, Skalierungsbedarfen und Team-Verantwortlichkeiten. \\
  Entscheidung & Jeder Bounded Context als eigenständiger Microservice mit eigenem Repository und eigener Datenbank. \\
  \rowcolor{tablegrey}
  Alternativen & (1)~\emph{Modularer Monolith} -- ernsthaft evaluiert:
  geringere Betriebskomplexität, einfachere lokale Entwicklung,
  ACID-Transaktionen über Modulgrenzen. Für ein 5--8-köpfiges Team und
  einen 6-Monats-MVP ist dies die objektiv einfachere Wahl.
  Abgelehnt wegen (a)~unterschiedlicher Skalierungsbedarfe:
  \texttt{registration} muss zur Spitzenzeit unabhängig skalieren;
  (b)~Erfahrung zeigt, dass Modulgrenzen im Monolithen mit der Zeit
  erodieren (``Big Ball of Mud''); (c)~Lernziel: Microservices
  realitätsnah dokumentieren.
  Fallback explizit geplant: falls Betriebsaufwand das Team überfordert,
  Rückbau auf modularen Monolithen (Risiko~R5).
  (2)~Zwei Services -- abgelehnt: fachliche Grenzen verwischen. \\
  Konsequenzen & (+)~Unabhängige Deployments. (+)~Team-Autonomie. (+)~Separate Skalierbarkeit. (-)~Erhöhte Betriebskomplexität. (-)~Eventual Consistency. \\
\end{tabularx}

% ── ADR-002 ──────────────────────────────────────────────────
\subsection{ADR-002: Asynchrone Kommunikation via Message Broker}

\begin{tabularx}{\textwidth}{L{3.5cm} X}
  \rowcolor{tableheader}
  \color{white}\textbf{Feld} & \color{white}\textbf{Inhalt} \\
  Status & Akzeptiert -- 2026-01-22 \\
  \rowcolor{tablegrey}
  Qualitätsziele & Reliability-Availability (primär), Maintainability, Functional Correctness \\
  Entscheidungstreiber & Resilienz bei Teilausfällen; Entkopplung von Sender und Empfänger; kein Kaskadenausfall durch synchrone Abhängigkeiten. \\
  \rowcolor{tablegrey}
  Entscheidung & Domain Events über RabbitMQ (Pub/Sub). \\
  Alternativen & (1)~Direkter REST-Aufruf -- abgelehnt: starke Kopplung, Kaskadenfehler. (2)~Polling der DB -- abgelehnt: ineffizient. (3)~Kafka -- zurückgestellt: Overkill für aktuelles Volumen. \\
  \rowcolor{tablegrey}
  Konsequenzen & (+)~Lose Kopplung. (+)~Resilienz. (-)~Eventual Consistency. (-)~Zusätzliche Infrastruktur. \\
\end{tabularx}

% ── ADR-003 ──────────────────────────────────────────────────
\subsection{ADR-003: Externer Identity Provider}

\begin{tabularx}{\textwidth}{L{3.5cm} X}
  \rowcolor{tableheader}
  \color{white}\textbf{Feld} & \color{white}\textbf{Inhalt} \\
  Status & Akzeptiert -- 2026-01-25 \\
  \rowcolor{tablegrey}
  Qualitätsziele & Security (primär), Maintainability \\
  Entscheidungstreiber & Kein eigenes Passwort-Management (Sicherheitsrisiko); MFA und SSO ohne Eigenentwicklung; DSGVO-konforme Datenhaltung beim IdP. \\
  \rowcolor{tablegrey}
  Entscheidung & Externer OIDC-konformer IdP (z.\,B.\ Keycloak). EduPlanner verwaltet keine Passwörter. \\
  Alternativen & (1)~Eigener Auth-Service -- abgelehnt: hohes Sicherheitsrisiko. (2)~Auth0/Okta -- in Evaluierung für v2.0. \\
  \rowcolor{tablegrey}
  Konsequenzen & (+)~Security-Best-Practices, MFA, SSO, DSGVO. (-)~Externe Abhängigkeit (R3). \\
\end{tabularx}

% ── ADR-004 ──────────────────────────────────────────────────
\subsection{ADR-004: Transactional Outbox für zuverlässige Events}

\begin{tabularx}{\textwidth}{L{3.5cm} X}
  \rowcolor{tableheader}
  \color{white}\textbf{Feld} & \color{white}\textbf{Inhalt} \\
  Status & Akzeptiert -- 2026-02-01 \\
  \rowcolor{tablegrey}
  Qualitätsziele & Functional Correctness (primär), Reliability-Availability \\
  Entscheidungstreiber & Atomizität zwischen DB-Schreibvorgang und Event-Publikation; Vermeidung von Dual-Write-Problem; Garantie: kein verlorenes Event. \\
  \rowcolor{tablegrey}
  Entscheidung & Events atomar mit DB-Änderung in \texttt{outbox}-Tabelle; Relay-Prozess publiziert asynchron (Technologie: ADR-009). \\
  Alternativen & (1)~Dual Write -- abgelehnt: kein Two-Phase-Commit. (2)~Event-Sourcing -- zu komplex. \\
  \rowcolor{tablegrey}
  Konsequenzen & (+)~Keine verlorenen Events. (+)~Atomare Konsistenz. (-)~Idempotente Handler erforderlich. \\
\end{tabularx}

% ── ADR-005 ──────────────────────────────────────────────────
\subsection{ADR-005: Angular SPA als Frontend-Technologie}

\begin{tabularx}{\textwidth}{L{3.5cm} X}
  \rowcolor{tableheader}
  \color{white}\textbf{Feld} & \color{white}\textbf{Inhalt} \\
  Status & Akzeptiert -- 2026-02-05 \\
  \rowcolor{tablegrey}
  Qualitätsziele & Usability (primär), Maintainability \\
  Entscheidungstreiber & Team-Kenntnisse in Vue.js und Angular vorhanden; SPA ermöglicht reaktive UX ohne Server-Round-Trip; TypeScript erhöht Wartbarkeit. \\
  \rowcolor{tablegrey}
  Entscheidung & Angular, TypeScript. \\
  Alternativen & (1)~Next.js -- zurückgestellt (v2.0). (2)~Vue.js -- als gleichwertige Alternative evaluiert; Angular gewählt aufgrund stärkerer Typisierung und integriertem Framework-Ökosystem. \\
  \rowcolor{tablegrey}
  Konsequenzen & (+)~Bekanntes Ökosystem. (+)~Starke Typisierung via TypeScript. (-)~Bundle-Grösse (Code Splitting erforderlich). (-)~SEO-Aufwand (kein SSR in v1.0). \\
\end{tabularx}

% ── ADR-006 ──────────────────────────────────────────────────
\subsection{ADR-006: PostgreSQL als primäre Datenbank}

\begin{tabularx}{\textwidth}{L{3.5cm} X}
  \rowcolor{tableheader}
  \color{white}\textbf{Feld} & \color{white}\textbf{Inhalt} \\
  Status & Akzeptiert -- 2026-02-08 \\
  \rowcolor{tablegrey}
  Qualitätsziele & Functional Correctness (primär), Reliability-Availability \\
  Entscheidungstreiber & ACID-Garantien für Transaktionen; JSONB-Support\footnote{JSONB
    (JSON Binary): speichert JSON-Daten binär dekomprimiert -- im Gegensatz zu JSON indizierbar
    (GIN-Index) und deutlich schneller abfragbar.} für flexible Payloads;
  Flyway-Kompatibilität; breite Java-Ökosystem-Unterstützung. \\
  \rowcolor{tablegrey}
  Entscheidung & PostgreSQL~15 (Managed) für alle Services. \\
  Alternativen & (1)~MySQL -- abgelehnt (JSONB-Support, Lizenz). (2)~MongoDB -- abgelehnt (ACID-Transaktionen). \\
  \rowcolor{tablegrey}
  Konsequenzen & (+)~Starke Konsistenz. (+)~JSONB-Flexibilität. (-)~Skalierbarkeit nur vertikal (Read-Replicas helfen). \\
\end{tabularx}

% ── ADR-007 ──────────────────────────────────────────────────
\subsection{ADR-007: Kein Service Mesh in v1.0}

Dieses ADR begründet den bewussten Verzicht auf ein Service Mesh (z.\,B.\ Linkerd oder Istio)
in der ersten Produktionsversion. Es korrespondiert direkt mit der technischen Schuld TS1 und
wird in v2.0 durch die Einführung von Linkerd abgelöst.

\begin{tabularx}{\textwidth}{L{3.5cm} X}
  \rowcolor{tableheader}
  \color{white}\textbf{Feld} & \color{white}\textbf{Inhalt} \\
  Status & Akzeptiert -- 2026-02-10 \\
  \rowcolor{tablegrey}
  Qualitätsziele & Maintainability (primär); Security (nachrangig) \\
  Entscheidungstreiber & Geringes Team-Know-how mit Service Mesh; erhöhter Betriebsaufwand für Sidecar-Infrastruktur übersteigt den Nutzen in v1.0; Budget- und Zeitrahmen des MVP. \\
  \rowcolor{tablegrey}
  Kontext & Ein Service Mesh ergänzt Kubernetes um mTLS\footnote{%
    \textit{mTLS (Mutual TLS):} Gegenseitige TLS-Authentifizierung, bei der sich beide Seiten
    (Client \textit{und} Server) mit einem Zertifikat ausweisen. Im Gegensatz zu einfachem TLS,
    das nur den Server authentifiziert, stellt mTLS sicher, dass auch der aufrufende Service
    vertrauenswürdig ist -- ein wichtiges Konzept zur Absicherung der Service-zu-Service-Kommunikation
    in Microservice-Architekturen.}
  zwischen Services, L7-Observability und automatisches Retry. Ohne Service Mesh kommunizieren
  Services intern über unverschlüsseltes HTTP innerhalb des privaten Kubernetes-Netzwerks. \\
  Entscheidung & Kein Service Mesh in v1.0. Netzwerk-Sicherheit wird durch Kubernetes Network Policies und TLS am API Gateway kompensiert. Tilgung: Linkerd in v2.0 (TS1). \\
  \rowcolor{tablegrey}
  Alternativen & (1)~Linkerd -- zurückgestellt auf v2.0: geringerer Overhead als Istio, bevorzugte Option. (2)~Istio -- abgelehnt: zu hohe operative Komplexität für Team-Grösse. \\
  Konsequenzen & (+)~Reduzierte Betriebskomplexität in v1.0. (+)~Schnellerer MVP-Launch. (-)~Kein mTLS intern (Sicherheitsrisiko TS1). (-)~Eingeschränkte L7-Observability. \\
\end{tabularx}

% ── ADR-008 ──────────────────────────────────────────────────
\subsection{ADR-008: Single-Tenant-Architektur in v1.0}

Dieses ADR begründet den bewussten Verzicht auf Multi-Tenancy in der ersten Version.
Es korrespondiert mit der technischen Schuld TS4 und wird in v2.0 durch ein
Schema-per-Tenant-Modell abgelöst (vgl.\ Kap.~14.3).

\begin{tabularx}{\textwidth}{L{3.5cm} X}
  \rowcolor{tableheader}
  \color{white}\textbf{Feld} & \color{white}\textbf{Inhalt} \\
  Status & Akzeptiert -- 2026-02-12 \\
  \rowcolor{tablegrey}
  Qualitätsziele & Maintainability (primär, kurzfristig); Functional Correctness \\
  Entscheidungstreiber & Zeitrahmen MVP (6 Monate); Multi-Tenancy erfordert durchgängige \texttt{tenant\_id}-Spalten, separates IdP-Realm und Tenant-Routing -- zu aufwändig für v1.0. Erste Zielgruppe: ein einzelner Kursanbieter (SkillBridge AG). \\
  \rowcolor{tablegrey}
  Kontext & Multi-Tenancy bezeichnet die Fähigkeit einer SaaS-Plattform, mehrere Mandanten (Kursanbieter) auf derselben Infrastruktur mit strikter Datentrennung zu betreiben. In v1.0 wird EduPlanner für einen einzigen Kursanbieter betrieben. \\
  Entscheidung & Single-Tenant-Betrieb in v1.0. Alle Services teilen eine gemeinsame Deployment-Instanz ohne mandantenbezogene Datentrennung. \\
  \rowcolor{tablegrey}
  Alternativen & (1)~Schema-per-Tenant von Anfang an -- abgelehnt: zu komplex für MVP-Zeitrahmen. (2)~Database-per-Tenant -- abgelehnt: zu hohe Infrastrukturkosten in v1.0. \\
  Konsequenzen & (+)~Deutlich vereinfachte Datenbankschicht in v1.0. (+)~Schnellerer MVP-Launch. (-)~Technische Schulden für v2.0 (TS4). (-)~Migration auf Multi-Tenancy erfordert Schema-Änderungen an allen Tabellen (R11). \\
\end{tabularx}

% ── ADR-009 ──────────────────────────────────────────────────
\subsection{ADR-009: Outbox Relay -- Technologiewahl}

Dieses ADR spezifiziert die konkrete Implementierung des Relay-Prozesses, der die
\texttt{outbox}-Tabelle aus ADR-004 ausliest und Events an RabbitMQ publiziert.

\begin{tabularx}{\textwidth}{L{3.5cm} X}
  \rowcolor{tableheader}
  \color{white}\textbf{Feld} & \color{white}\textbf{Inhalt} \\
  Status & Akzeptiert -- 2026-02-15 \\
  \rowcolor{tablegrey}
  Qualitätsziele & Functional Correctness (primär), Reliability-Availability \\
  Entscheidungstreiber & Einfache, wartbare Lösung ohne zusätzliche externe Abhängigkeiten; muss \textit{at-least-once}-Delivery garantieren; muss idempotent konsumierbar sein. \\
  \rowcolor{tablegrey}
  Kontext & Der Relay-Prozess ist ein Background-Thread innerhalb jedes Services, der periodisch die \texttt{outbox}-Tabelle nach \texttt{PENDING}-Einträgen abfragt, diese an RabbitMQ sendet und den Status auf \texttt{SENT} setzt. Bei Fehlern nach 5 Versuchen landet das Event in der Dead-Letter-Queue (DLQ). \\
  Entscheidung & Polling-basierter Relay-Thread (Spring \texttt{@Scheduled}, Intervall 500\,ms) innerhalb jedes Services. Kein separater Prozess, keine externe Abhängigkeit. \\
  \rowcolor{tablegrey}
  Alternativen & (1)~Debezium\footnote{\textit{Debezium}: Open-Source-Plattform für Change Data Capture (CDC), die Datenbankänderungen aus dem Transaction Log ausliest und als Events publiziert. Leistungsfähiger als Polling, aber deutlich komplexer im Betrieb.} (CDC) -- zurückgestellt auf v2.0: leistungsfähiger, aber höhere operative Komplexität. (2)~Separater Relay-Service -- abgelehnt: unnötige Infrastruktur für aktuelles Volumen. \\
  Konsequenzen & (+)~Keine zusätzliche Infrastruktur. (+)~Einfach zu debuggen. (-)~Polling-Latenz ($\leq$\,500\,ms). (-)~Erhöhte DB-Last bei hohem Event-Volumen. \\
\end{tabularx}

% ── ADR-010 ──────────────────────────────────────────────────
\subsection{ADR-010: Cloud-Provider -- Microsoft Azure}

Dieses ADR dokumentiert die Wahl von Microsoft Azure als Cloud-Provider für alle
Produktions- und Staging-Umgebungen von EduPlanner (vgl.\ Kap.~7.1.1).

\begin{tabularx}{\textwidth}{L{3.5cm} X}
  \rowcolor{tableheader}
  \color{white}\textbf{Feld} & \color{white}\textbf{Inhalt} \\
  Status & Akzeptiert -- 2026-02-20 \\
  \rowcolor{tablegrey}
  Qualitätsziele & Reliability-Availability (primär), Security, Maintainability \\
  Entscheidungstreiber & DSGVO/DSG-Compliance mit Datenresidenz Schweiz; Managed Services für PostgreSQL und Kubernetes; bestehendes Team-Know-how mit Azure; Budget $\leq$ CHF~800/Monat. \\
  \rowcolor{tablegrey}
  Kontext & EduPlanner muss personenbezogene Daten ausschliesslich in der EU bzw.\ der Schweiz speichern (Kap.~2.3). Nicht alle Cloud-Provider bieten Schweizer Regionen mit vollständigem Managed-Service-Angebot an. \\
  Entscheidung & Microsoft Azure mit primärer Region \textit{Switzerland North} (Zürich) und Failover-Region \textit{Switzerland West} (Genf). Eingesetzte Managed Services: Azure Kubernetes Service (AKS~1.28+), Azure Database for PostgreSQL~15 Flexible Server, Azure Key Vault, Azure Monitor. \\
  \rowcolor{tablegrey}
  Alternativen & (1)~AWS -- abgelehnt: Schweizer Region (\textit{eu-central-2}) verfügbar, aber geringeres Team-Know-how; höhere Migrationskosten. (2)~Google Cloud -- abgelehnt: kein Schweizer Rechenzentrum zum Entscheidungszeitpunkt. (3)~On-Premise -- ausgeschlossen durch Randbedingung (Kap.~2.1). \\
  Konsequenzen & (+)~Datenresidenz Schweiz erfüllt DSGVO/DSG. (+)~Managed PostgreSQL: automatisches Failover, PITR, integriertes Monitoring. (+)~Team-Know-how reduziert Einarbeitungsaufwand. (-)~Gewisser Vendor Lock-in (mitigiert durch cloud-agnostische Tools: K8s, PostgreSQL, RabbitMQ -- vgl.\ R7). \\
\end{tabularx}

% ── ADR-011 ──────────────────────────────────────────────────
\subsection{ADR-011: Contract-First mit OpenAPI~3.1 und Swagger~UI}

Dieses ADR begründet den Contract-First-Ansatz mit OpenAPI~3.1.0 als
verbindliche Entwicklungsstrategie für alle REST-Schnittstellen von
EduPlanner sowie die Bereitstellung der interaktiven Swagger~UI pro Service.
Die konsolidierte Spezifikation liegt in Anhang~D.

\begin{tabularx}{\textwidth}{L{3.5cm} X}
  \rowcolor{tableheader}
  \color{white}\textbf{Feld} & \color{white}\textbf{Inhalt} \\
  Status & Akzeptiert -- 2026-02-22 \\
  \rowcolor{tablegrey}
  Qualitätsziele & Maintainability (primär), Functional Correctness, Usability \\
  Entscheidungstreiber &
  API-First als verbindliche Randbedingung (Kap.~2.1);
  vollständige OpenAPI-Dokumentation als Go-Live-Kriterium v1.0 (Kap.~1.1);
  Contract Testing (TS2) setzt maschinenlesbare Spec voraus;
  parallele Frontend-/Backend-Entwicklung ohne Warten auf Implementierung. \\
  \rowcolor{tablegrey}
  Kontext &
  Im \emph{Code-First}-Ansatz (z.\,B.\ Springdoc-Annotationen) entsteht
  die Spec nachgelagert aus dem Code -- sie beschreibt, was implementiert
  wurde, nicht was vereinbart wurde. Dies führt zu Drift zwischen Spec
  und Implementierung sowie zu blockierter Frontend-Entwicklung.
  Im \emph{Contract-First}-Ansatz ist die YAML-Datei die einzige
  Quelle der Wahrheit: Sie entsteht zuerst, wird in Git versioniert,
  und Server-Stubs werden daraus generiert. Der Compiler erzwingt
  Konformität. \\
  Entscheidung &
  \textbf{Contract-First} mit OpenAPI~3.1.0 (YAML).
  Die Datei \texttt{openapi/}\allowbreak \texttt{eduplanner-}\allowbreak \texttt{api.yaml} wird manuell gepflegt.
  Server-Stubs (Spring-Boot-Interfaces) werden via
  \texttt{openapi-generator-}\allowbreak\texttt{maven-plugin} generiert
  (\texttt{interfaceOnly=true},\newline\texttt{useSpringBoot3=true}).
  Das Dev-Team implementiert ausschliesslich die generierten Interfaces.
  Jeder Service stellt die interaktive Swagger~UI unter
  \texttt{/swagger-ui/}\allowbreak\texttt{index.html} bereit
  (\texttt{springdoc-openapi-}\allowbreak\texttt{starter-webmvc-ui});
  in Produktion nur aus dem internen Netz erreichbar.
  Die Spec wird in der CI-Pipeline mit \texttt{vacuum} (Linting) und
  \texttt{openapi-generator} (Generierbarkeit) validiert --
  ein ungültiger Vertrag blockiert den Build. \\
  \rowcolor{tablegrey}
  Alternativen &
  (1)~\emph{Code-First / Springdoc:} abgelehnt -- Spec entsteht
  nachgelagert, Frontend-Entwicklung blockiert, Drift möglich.
  (2)~\emph{GraphQL:} zurückgestellt auf v2.0 als optionale Schicht
  neben REST (Kap.~14.3).
  (3)~\emph{AsyncAPI:} evaluiert für Domain Events (RabbitMQ);
  zurückgestellt auf v1.1 -- ergänzt OpenAPI, ersetzt es nicht.
  (4)~\emph{Keine formale Spec:} abgelehnt -- Contract Testing (TS2)
  und externes Onboarding nicht möglich. \\
  Konsequenzen &
  (+)~Spec und Implementierung sind strukturell nie aus dem Sync
  (Compiler-Enforcement). \newline
  (+)~Frontend kann mit generierten TypeScript-Clients parallel
  entwickeln. \newline
  (+)~Swagger~UI (\texttt{/swagger-ui/index.html}) bietet
  interaktives API-Testing ohne externes Tool. \newline
  (+)~Basis für Pact Contract Tests (TS2-Tilgung Sprint~8). \newline
  (-)~Initiale Spec-Erstellung erfordert Disziplin; kein ``quick
  and dirty'' Endpoint möglich. \newline
  (-)~OpenAPI beschreibt nur synchrone REST -- für asynchrone
  Domain Events ist AsyncAPI ergänzend nötig (v1.1). \\
\end{tabularx}
